%% LaTeX2e class for student theses
%% sections/abstract_de.tex
%% 
%% Karlsruhe Institute of Technology
%% Institute for Program Structures and Data Organization
%% Chair for Software Design and Quality (SDQ)
%%
%% Dr.-Ing. Erik Burger
%% burger@kit.edu
%%
%% Version 1.6, 2024-06-07

\Abstract
Die Effizienz der Energiespeicherung über einen Metallträger hängt maßgeblich sowohl von der Effizienz der Reduktion von Eisenoxid zu Eisen als auch von der Oxidation des Eisens zur Energiefreisetzung ab. Aktuelle werden Experimente zum Einsatz eines Wirbelschichtreaktors für die Reaktion eingesetzt. Die Strömungsdynamik, die Partikelbewegung und die Partikel-Partikel-Wechselwirkungen innerhalb einer Wirbelschicht sind komplex. Ein Ansatz, um das Strömungs- und Bewegungsverhalten im Reaktor besser zu verstehen, ist die Nutzung numerischer Simulationen.
Um die Rechenzeit zu reduzieren, werden Partikel-Partikel-Wechselwirkungen auf Grundlage von Eigenschaften modelliert, die auf einem Eulerschen Gitter berechnet werden, anstatt vollständig aufgelöst zu werden. Dies erfolgt mittels der \MPns-Methode. Wir betrachten die \MPns-Methode genauer in der Form, in der sie ursprünglich vorgestellt wurde, einschließlich der verwendeten Submodelle. Dazu gehört insbesondere das Pack-Modell, das eine Überpackung verhindert, indem es Partikel beschleunigt, sobald sie sich nahe einer dichtesten Packung befinden.
Die Diskretisierung basiert auf einem einfachen Finite-Differenzen-Schema, das auf die Erhaltungsgleichungen der Fluid- und Partikelphase angewendet wird. Wir formulierten die verwendeten Partikeleigenschaften so um, dass sie ausschließlich von den Werten des Eulerschen Gitters abhängen. Daraus ergeben sich vier Gleichungen: die diskretisierten Kontinuitäts- und Impulsgleichungen für sowohl die Fluid- als auch die Feststoffphase, die mithilfe vorläufig bestimmter Partikelpositionen gelöst wurden. Dies unterscheidet sich von der Implementierung in \OF.
Die Partikel werden explizit verfolgt und ihre Bewegung basiert auf den Newtonschen Bewegungsgleichungen. Mehrere Submodelle sind in \OF implementiert. Hier diskutieren wir das Isotropiemodell, das explizite Harris-und-Crighton-Pack-Modell, zwei Strömungswiderstands-Modelle sowie den \lstinline|PIMPLE|-Algorithmus, welcher die für die Submodelle benötigten Größen auf dem Eulerschen Gitter löst.
Wir führen eine Gitterstudie durch und schätzen die Konvergenzordnung ab. Zur Validierung des Solvers verarbeiten wir Daten der Arbeitsgruppe "Verbrennungstechnik". In einer parametrischen Studie vergleichen wir die resultierende Betthöhe und den Druckverlust mit simulierten Ergebnissen für das Harris‑und‑Crighton‑Modell in Kombination mit unterschiedlichen Strömungswiderstands-Modellen und einem Isotropiemodell. Wir zeigen, dass die \MP-Methode bei geeigneter Parameterauswahl die Betthöhe gut vorhersagen kann, jedoch vom experimentell ermittelten Druckverlust abweicht.
